\chapter*{About These Notes}
\addcontentsline{toc}{chapter}{About These Notes}

\paragraph{Goal.} These notes accompany a short course for readers who know real
analysis and mathematical statistics but have not met queueing theory. They build,
from definitions, the queueing tools needed to model one architecture of large
language model (LLM) serving, \emph{prefill--decode disaggregation}, and to prove the
results a scheduler for it needs. Measure-theoretic probability is not needed; we use
random variables, expectation, conditional expectation, the strong law of large
numbers and power series freely.

\paragraph{The system.} An LLM request is served in two phases. \emph{Prefill} reads
the whole prompt in one pass and produces the first token; the time from the
request's arrival to its first token is the \emph{time to first token} (TTFT).
\emph{Decode} then produces one token per step until the answer is complete. Both
phases need the attention state of every prompt token, the \emph{KV cache}. In a
disaggregated deployment the two phases run on different machines:
\begin{itemize}
\item \emph{prefill instances} run prefills only, one after another; they keep a
  \emph{prefix cache} of recent prompts, so that a request that starts with a cached
  prefix (a \emph{hit}) computes only its new tokens, while one whose prefix was
  evicted (a \emph{miss}) recomputes its whole context;
\item a \emph{link} carries each request's KV cache from its prefill instance to a
  decode instance;
\item \emph{decode instances} run decodes only, advancing all their running requests
  by one token per step.
\end{itemize}
Both kinds of instance hold KV in a memory pool of fixed size, so a request is admitted
to an instance only when its context fits. An \emph{agent} is a session of many such
requests (\emph{turns}): after each turn it calls a tool, pauses, and sends the next
turn, whose prompt is the previous one with the tool's output appended. Requests
arrive at random, wait for memory and for servers, receive service and leave: this is
the object of queueing theory.

\paragraph{Why disaggregation makes the queueing model clean.} When prefill and decode
share a machine, each slows the other and the queue each one sees depends on the
other's load. Separated, each instance is a single, recognisable queue: a first-in
first-out server for prefill, a processor-sharing batch for decode, a shared-bandwidth
link in between, and a multi-server memory station in front of each instance. The
notes study these pieces one at a time and then put them together.

\paragraph{Map.} What the course proves and where:

\begin{center}\small
\begin{tabular}{@{}p{0.34\linewidth}p{0.6\linewidth}@{}}
\toprule
Question & Tools and results \\
\midrule
How long does a request wait for memory? &
  birth--death processes (\cref{thm:birthdeath}), M/M/$m$ and the Erlang C formula
  (\cref{sec:L2}) \\
Does the decode batch care about the length distribution? &
  processor sharing and insensitivity (\cref{thm:insens}), \cref{thm:L-mono} \\
What does one prefix miss cost the prefill instance? &
  PASTA (\cref{thm:pasta}), residual service time (\cref{lem:residual}),
  Pollaczek--Khinchine (\cref{thm:pk}), the price of a miss (\cref{thm:price}) \\
What does it cost the link and the decode instance? &
  the price at a processor-sharing station (\cref{prop:decode-price}) \\
How many sessions can be live? &
  closed systems and mean value analysis (\cref{prop:finite}), Campbell's theorem
  (\cref{thm:campbell}) \\
What should each instance keep, and for how long? &
  covering knapsack and the threshold rule (\cref{thm:threshold}),
  price-blind keys (\cref{prop:blind}), \cref{sec:L5} \\
Can misses feed themselves? &
  monotone fixed points (\cref{sec:L5}) \\
\bottomrule
\end{tabular}
\end{center}

\paragraph{Plan.} \Cref{sec:L1} covers arrival processes, Little's law and the
queueing model of the disaggregated system; \cref{sec:L2} Markovian queues,
memory slots and processor sharing; \cref{sec:L3} the prefill instance as an M/G/1
queue and the price of a miss; \cref{sec:L4} closed systems, memory occupancy and
eviction; and \cref{sec:L5} assembles the pieces into a model of the disaggregated
system with its transfer overhead, prices and feedback. The exercises at the end of
each lecture are meant to be solved before the next one. Good introductory textbooks
are \cite{kleinrock1975,harchol2013}; \cite{asmussen2003} is the rigorous reference.
